\chapter{Batching and Serving Engines}\label{ch:serving}

\Cref{ch:inference} ended with a sobering table: a single decoding stream uses well under 1\% of a
modern GPU's arithmetic, because every token drags the entire model through the memory system.
The only way to amortize that cost is to decode many sequences at once. This chapter is about
doing so well. We start where every practitioner starts: pad a few prompts into a rectangle and
run them together. That simple idea already speeds up every workload of Part~VI (evaluation,
self-consistency sampling, GRPO rollouts and distillation), and we walk through each one. Then we
take the serving-system view. Requests arrive at random times with random lengths; rectangles waste
compute on padding and memory on reservations, and fast requests get stuck behind slow ones. The
fixes, iteration-level (continuous) batching, paged KV memory, prefix sharing, chunked prefill, and
disaggregation of prefill from decode, form the core of every modern engine (vLLM, SGLang,
TensorRT-LLM, TGI). We finish by building a working miniature of such an engine,
\ident{llmsys.serving}, that produces exactly the same tokens as our reference generator while
batching, paging, sharing and preempting underneath.

\begin{figure}[t]
\centering
\begin{tikzpicture}[
  stage/.style={draw=sysblue,fill=sysblue!8,rounded corners=2pt,minimum width=2.45cm,
    minimum height=1.0cm,align=center,font=\small},
  sys/.style={stage,draw=napkin,fill=napkin!10},
  num/.style={circle,fill=sysblue,text=white,font=\scriptsize\bfseries,inner sep=1.2pt},
  arr/.style={-{Stealth[length=2mm]},thick,sysblue!70}]
\node[stage] (a) at (0,0) {Why batching\\helps};
\node[stage] (b) at (2.9,0) {Padding \&\\masks};
\node[stage] (c) at (5.8,0) {Batched eval,\\sampling, RL, KD};
\node[stage] (d) at (8.7,0) {Single vs.\\batch trade-offs};
\node[sys] (e) at (11.6,0) {Continuous\\batching};
\node[sys] (f) at (11.6,-1.9) {Paged KV\\memory};
\node[sys] (g) at (8.7,-1.9) {Prefix sharing};
\node[sys] (h) at (5.8,-1.9) {Chunked prefill\\\& scheduling};
\node[sys] (i) at (2.9,-1.9) {Disaggregation};
\node[sys,draw=labpurple,fill=labpurple!10] (j) at (0,-1.9) {Mini engine\\+ Lab 19};
\foreach \x/\y in {a/b,b/c,c/d,d/e,e/f,f/g,g/h,h/i,i/j} \draw[arr] (\x) -- (\y);
\foreach \n/\k in {a/1,b/2,c/3,d/4,e/5,f/6,g/7,h/8,i/9,j/10} \node[num] at (\n.north west) {\k};
\node[font=\scriptsize\itshape,text=sysblue] at (5.8,0.85) {batching in our own scripts (offline)};
\node[font=\scriptsize\itshape,text=napkin] at (5.8,-2.75) {batching in a server (online, random arrivals)};
\end{tikzpicture}
\caption{Roadmap. The top row batches known workloads offline; the bottom row designs an online
engine for a stream of unpredictable requests.}
\label{fig:serving:roadmap}
\end{figure}

% =====================================================================
\section{Latency, throughput, and why batching helps}\label{sec:serving:why}

Two performance goals pull in different directions:
\begin{itemize}
  \item \textbf{Latency}: how quickly \emph{one} prompt gets its answer.
  \item \textbf{Throughput}: how many prompts (or tokens) are processed per unit time.
\end{itemize}
Single-sequence generation, as in \cref{ch:inference}, minimizes latency for one user, is easiest
to read and debug, and is what we use to explain algorithms throughout the book. Batched generation
stacks several sequences into one tensor and advances all of them with each forward pass
(\cref{fig:serving:singlebatch}). It is a \emph{throughput} optimization.

\begin{figure}[ht]
\centering
\begin{tikzpicture}[
  q/.style={draw=fdata,fill=fdata!12,rounded corners=1pt,minimum width=2.3cm,minimum height=0.48cm,font=\scriptsize},
  a/.style={draw=napkin,fill=napkin!15,rounded corners=1pt,minimum height=0.48cm,font=\scriptsize},
  llm/.style={draw=fcompute,fill=fcompute!15,rounded corners=2pt,minimum width=1.0cm,minimum height=0.5cm,font=\scriptsize\bfseries},
  arr/.style={-{Stealth[length=1.4mm]},thick}]
\node[font=\small\bfseries,anchor=west] at (-1.3,1.0) {sequential: one prompt per pass};
\foreach \i/\p/\r in {0/{Solve $3x-1=8$}/{$x=3$},1/{Compute $9^2-4^2$}/{$65$},2/{Simplify $18/6+2$}/{$5$}} {
  \node[q] (Q\i) at (4.1*\i,0.25) {\p};
  \node[llm] (L\i) at (4.1*\i,-0.5) {LLM};
  \node[a,minimum width=1.2cm] (A\i) at (4.1*\i,-1.25) {\r};
  \draw[arr] (Q\i) -- (L\i); \draw[arr] (L\i) -- (A\i);
  \node[circle,fill=sysblue,text=white,font=\tiny\bfseries,inner sep=1pt] at (4.1*\i-1.35,0.25) {\the\numexpr\i+1\relax};
}
\draw[-{Stealth},black!50] (-1.3,-1.8) -- (10.0,-1.8) node[right,font=\scriptsize]{time};
\begin{scope}[yshift=-3.3cm]
\node[font=\small\bfseries,anchor=west] at (-1.3,1.0) {batched: all prompts in one pass};
\node[q] at (0,0.45) {Solve $3x-1=8$};
\node[q] at (0,0.0) {Compute $9^2-4^2$};
\node[q] at (0,-0.45) {Simplify $18/6+2$};
\node[llm,minimum height=1.4cm,minimum width=1.4cm] (LB) at (2.6,0) {LLM};
\node[a,minimum width=1.2cm] at (4.6,0.45) {$x=3$};
\node[a,minimum width=1.2cm] at (4.6,0.0) {$65$};
\node[a,minimum width=1.2cm] at (4.6,-0.45) {$5$};
\draw[arr] (1.2,0) -- (LB); \draw[arr] (LB) -- (3.95,0);
\node[font=\scriptsize,align=left,anchor=west,text=black!70] at (5.6,0)
  {one weight read serves 3 rows;\\the tensor is $(3,T)$ instead of $(1,T)$};
\draw[-{Stealth},black!50] (-1.3,-0.95) -- (10.0,-0.95) node[right,font=\scriptsize]{time};
\end{scope}
\end{tikzpicture}
\caption{Sequential versus batched generation. In the batched case the prompts share every
forward pass, so each weight read from memory is amortized over all rows.}
\label{fig:serving:singlebatch}
\end{figure}

The reason batching works is the roofline argument of \cref{sec:inference:roofline}: the weight
GEMMs of a decode step have arithmetic intensity $\approx B$ (in bf16), so on a GPU with a ridge point
in the hundreds, raising $B$ from 1 to 64 multiplies useful FLOPs per byte read by 64 while the step
time barely changes. But batching is not free, and it is not always faster:
\begin{itemize}
  \item Prompts have different lengths, and tensors are rectangles: shorter rows are padded, and
    the padding costs compute and memory without producing anything.
  \item Rows finish at different times, so a naive batch runs until its longest row is done.
  \item Each row has its own KV cache, so batch size is bounded by memory, and attention over the
    cache does not get cheaper per row as the batch grows.
  \item On CPUs or for tiny models, the step is not dominated by weight reads in the first place,
    and the bookkeeping (masks, padding) can eat the gain.
\end{itemize}

\begin{keyidea}
Batching is a throughput optimization with a latency cost. Whether it pays off is decided by
where the step time goes: amortizable work (weight reads, kernel launches) versus per-row work
(KV reads, attention, padding).
\end{keyidea}

% =====================================================================
\section{Running batched generation}\label{sec:serving:batched}

The core obstacle is that prompts tokenize to different lengths. To run them together we pad
shorter rows with a placeholder id up to the batch's maximum length, and keep a boolean mask that
records which positions are real. Which id we pad with does not matter, because padded positions
will be masked out; reusing the end-of-text token is common.

Where should the padding go? For \emph{generation} we pad on the \textbf{left}. Left padding aligns
the last real token of every prompt in the final column, so after the prefill the next-token logits
of all rows sit at position $-1$, and each decode step appends one column to everyone
(\cref{fig:serving:masks}). With right padding, every row's ``last real token'' would be at a
different index, and new tokens would have to be written into the middle of the tensor.

\begin{lstlisting}[caption={Left-padded batched generation with the core helper.},label={lst:serving:batch}]
from llmsys.generate import generate_batch

problems = ["2+2?", "3+3=6?", "Is 91 prime?"]
prompts = [tok.encode(render_prompt(p)) for p in problems]   # ragged lists
outs = generate_batch(model, prompts, max_new_tokens=64,
                      pad_id=tok.pad_token_id, eos_id=tok.pad_token_id)
for o in outs:                       # each row trimmed at its own EOS
    print(tok.decode(o))
\end{lstlisting}

Inside \ident{generate_batch} the steps are: build the left-padded $(B,T)$ id tensor and a
$(B,T+n)$ validity mask; run one prefill over the whole rectangle with a KV cache; then for each
step sample one token per row, append it, and run a decode step with the mask extended by one
always-valid column.

\paragraph{Finished rows.} When a row emits EOS, three strategies are possible
(\cref{fig:serving:eos}): (1) process prompts one by one; (2) keep the finished row in the batch,
force it to keep emitting EOS so shapes stay aligned, and trim afterwards; (3) remove finished
rows from the active batch, which saves compute but requires slicing the KV cache of every layer
(and re-indexing masks) each time a row leaves. Strategy (2) is simplest and is what
\ident{generate_batch} does (it stops recording tokens for finished rows). Strategy (3) helps when
output lengths are very skewed, and it is the seed of continuous batching (\cref{sec:serving:continuous}):
once rows can \emph{leave} a running batch, it is natural to let new rows \emph{join} it.

\begin{figure}[ht]
\centering
\begin{tikzpicture}[
  c/.style={draw=black!30,minimum width=0.85cm,minimum height=0.46cm,font=\scriptsize,inner sep=0pt},
  e/.style={c,fill=warnred!20},
  w/.style={c,fill=black!6,text=black!35},
  lab/.style={font=\scriptsize,anchor=east}]
\foreach \s in {1,...,4} \node[font=\scriptsize\bfseries] at (2.0+0.9*\s,0.55) {step \s};
% (1) sequential
\node[font=\scriptsize\bfseries,anchor=west] at (-1.9,0.55) {(1) one at a time};
\node[lab] at (2.3,0) {$3x-1=8$};
\foreach \s/\t in {1/$x$,2/$=$,3/$3$,4/eos} {\ifnum\s=4 \node[e] at (2.0+0.9*\s,0) {\t}; \else \node[c] at (2.0+0.9*\s,0) {\t}; \fi}
% (2) keep finished rows
\begin{scope}[yshift=-1.0cm]
\node[font=\scriptsize\bfseries,anchor=west] at (-1.9,0.4) {(2) batch, keep finished rows};
\foreach \r/\p/\a/\b/\cc/\d in {0/{$3x-1=8$}/{$x$}/{$=$}/{$3$}/eos,1/{$9^2-4^2$}/{65}/eos/eos/eos,2/{$18/6+2$}/{5}/{.}/eos/eos} {
  \node[lab] at (2.3,-0.5*\r) {\p};
  \foreach \s/\t in {1/\a,2/\b,3/\cc,4/\d} {
    \node[c] (n) at (2.0+0.9*\s,-0.5*\r) {\t};
  }
}
\foreach \x/\y in {4/0,2/-0.5,3/-0.5,4/-0.5,3/-1.0,4/-1.0} \node[e] at (2.0+0.9*\x,\y) {eos};
\foreach \x/\y in {3/-0.5,4/-0.5,4/-1.0} \node[w] at (2.0+0.9*\x,\y) {eos};
\node[font=\scriptsize,anchor=west,text=black!60,align=left] at (6.1,-0.5) {grey: wasted compute\\on already-finished rows};
\end{scope}
% (3) drop finished
\begin{scope}[yshift=-3.0cm]
\node[font=\scriptsize\bfseries,anchor=west] at (-1.9,0.4) {(3) batch, drop finished rows};
\foreach \r/\p in {0/{$3x-1=8$},1/{$9^2-4^2$},2/{$18/6+2$}} \node[lab] at (2.3,-0.5*\r) {\p};
\node[c] at (2.9,0) {$x$}; \node[c] at (3.8,0) {$=$}; \node[c] at (4.7,0) {$3$}; \node[e] at (5.6,0) {eos};
\node[c] at (2.9,-0.5) {65}; \node[e] at (3.8,-0.5) {eos};
\node[c] at (2.9,-1.0) {5}; \node[c] at (3.8,-1.0) {.}; \node[e] at (4.7,-1.0) {eos};
\node[font=\scriptsize,anchor=west,text=black!60,align=left] at (6.1,-0.5) {batch shrinks 3 $\to$ 2 $\to$ 1;\\KV cache rows must be sliced};
\end{scope}
\end{tikzpicture}
\caption{Three ways to handle sequences that finish at different steps.}
\label{fig:serving:eos}
\end{figure}

% =====================================================================
\section{Padding and attention masks}\label{sec:serving:masks}

For a single unpadded prompt, the model builds the standard causal mask: position $i$ may attend
to positions $j\le i$. With padding we must additionally make sure no real token ever attends to a
pad token, otherwise the pad embeddings leak into the computation and the batched output differs
from the single-sequence output. The combined rule is
\begin{equation}\label{eq:serving:mask}
  \text{forbidden}(b,i,j) \;=\; \underbrace{[\,j>i\,]}_{\text{causal}} \;\lor\; \underbrace{[\,\neg\text{valid}(b,j)\,]}_{\text{key is padding}},
\end{equation}
which is exactly what \ident{llmsys.layers.build_mask} computes from a $(B,T_k)$ validity mask
(\cref{fig:serving:masks}). Two details matter in practice:
\begin{itemize}
  \item \textbf{Padded queries.} A pad position on the left has no valid key at all (its only
    allowed keys, those $\le i$, are all padding), so its softmax row is $\softmax(-\infty,\ldots)$ =
    NaN. These rows produce garbage that no real token ever reads, but NaNs must not propagate
    (e.g.\ through a later normalization or a $0\cdot\text{NaN}$ in a fused kernel). Our explicit attention
    zeroes them with \code{nan\_to\_num}; production kernels special-case fully-masked rows.
  \item \textbf{Decode steps extend the mask.} Each new column is valid for every row, so the
    validity mask grows by a column of \code{True} per step, while the left padding stays masked
    forever.
\end{itemize}

\begin{figure}[ht]
\centering
\begin{tikzpicture}[m/.style={minimum size=0.48cm,draw=black!25,font=\tiny,inner sep=0pt}]
\newcommand{\drawmask}[4]{% x offset, pad count, title, label
  \begin{scope}[xshift=#1cm]
  \node[font=\scriptsize\bfseries] at (1.2,0.75) {#3};
  \foreach \i in {0,...,5} { \foreach \j in {0,...,5} {
    \pgfmathtruncatemacro{\pad}{(\j<#2) ? 1 : 0}
    \pgfmathtruncatemacro{\qpad}{(\i<#2) ? 1 : 0}
    \pgfmathtruncatemacro{\fut}{(\j>\i) ? 1 : 0}
    \ifnum\fut=1 \node[m,fill=black!45] at (0.48*\j,-0.48*\i) {};
    \else\ifnum\pad=1 \node[m,fill=warnred!45] at (0.48*\j,-0.48*\i) {};
    \else \node[m,fill=keygreen!35] at (0.48*\j,-0.48*\i) {}; \fi\fi
    \ifnum\qpad=1 \ifnum\fut=0 \node[m,pattern=north east lines,pattern color=warnred] at (0.48*\j,-0.48*\i) {}; \fi\fi
  }}
  \node[font=\tiny,align=center] at (1.2,-3.1) {#4};
  \end{scope}}
\drawmask{0}{0}{causal only}{6 real tokens}
\drawmask{4.0}{2}{causal + key padding}{2 left pads, 4 real tokens}
\begin{scope}[xshift=8.6cm,yshift=-0.6cm]
\node[m,fill=keygreen!35] at (0,0) {}; \node[font=\scriptsize,anchor=west] at (0.35,0) {may attend};
\node[m,fill=black!45] at (0,-0.6) {}; \node[font=\scriptsize,anchor=west] at (0.35,-0.6) {future (causal)};
\node[m,fill=warnred!45] at (0,-1.2) {}; \node[font=\scriptsize,anchor=west] at (0.35,-1.2) {padding key};
\node[m,pattern=north east lines,pattern color=warnred] at (0,-1.8) {}; \node[font=\scriptsize,anchor=west,align=left] at (0.35,-1.95) {padding query\\(row is all-masked)};
\end{scope}
\node[font=\scriptsize] at (-0.55,-1.2) {\rotatebox{90}{query $i$}};
\node[font=\scriptsize] at (1.2,0.38) {key $j$};
\end{tikzpicture}
\caption{Masks for a 6-column batch row. Left: a row without padding. Right: the same layout with two
left pads; real queries attend only to real, non-future keys, which reproduces the 4-token causal mask
of the unpadded prompt in the lower-right corner.}
\label{fig:serving:masks}
\end{figure}

\paragraph{Positions under left padding.} Masking makes pad \emph{keys} invisible, but the real tokens
of a padded row now sit at column indices $p, p+1,\ldots$ instead of $0,1,\ldots$, where $p$ is the
row's pad count. Whether that shift matters depends on the positional scheme.

\begin{proposition}[RoPE is invariant to a uniform shift]\label{prop:serving:rope}
Let $R_m$ be the RoPE rotation for position $m$. For any query at position $m$ and key at position
$n$, and any offset $p$,
$\langle R_{m+p}\,q,\ R_{n+p}\,k\rangle = \langle R_m\,q,\ R_n\,k\rangle$.
\end{proposition}
\begin{proof}
Each 2-D frequency pair is rotated by an angle proportional to the position, so
$R_m^\top R_n = R_{n-m}$ (rotations compose additively and $R_m^\top=R_{-m}$). Hence
$\langle R_{m+p}q,R_{n+p}k\rangle = q^\top R_{m+p}^\top R_{n+p}k = q^\top R_{n-m}k$, which does not depend on $p$.
\end{proof}

So for Qwen3, left padding changes no attention score of a real token, and batched outputs equal
single-sequence outputs up to floating-point reordering (a test in the core
test suite, \ident{tests/test_core.py}, checks exactly this). GPT-2 instead adds a
\emph{learned absolute} position embedding, so a row shifted by $p$ sees different embeddings and
its outputs change slightly. The fix is to pass explicit \emph{position ids} per row,
$\text{pos}(b,t)=\max(0,\,t-p_b)$, which is what Hugging Face's \code{position\_ids} argument and
every production engine do. (Our GPT-2 class takes positions from the cache length; extending it to
per-row positions is \cref{ex:serving:positions}.)

\begin{pitfall}[Right padding for training, left padding for generation]
Training and scoring (classification, distillation, log-probability computation) usually pad on the
\emph{right}, because labels are aligned with inputs and nothing is generated after the last token;
loss masks (label $-100$, \cref{ch:instruct}) remove pad positions. Generation pads on the \emph{left}.
Mixing them up is the most common batched-generation bug: with right padding, \code{logits[:, -1]} reads
the prediction after a \emph{pad} token for every row but the longest.
\end{pitfall}

% =====================================================================
\section{Batching across the reasoning workflow}\label{sec:serving:workflows}

Every compute-heavy step of Part~VI benefits from batching, but the steps differ in an important way:
some batch \emph{different} prompts of different lengths (padding needed), others batch many samples
of the \emph{same} prompt (no padding, and a shared prefix that a smart engine computes only once).
\Cref{tab:serving:workflows} summarizes; the subsections give the details.

\begin{table}[ht]
\centering\small
\caption{What gets batched in each workflow of the book.}
\label{tab:serving:workflows}
\begin{tabularx}{\linewidth}{lXll}
\toprule
Workflow & Rows of the batch & Padding? & Shared prefix?\\
\midrule
MATH-500 evaluation (\cref{ch:eval}) & different problems & yes (left) & only the template\\
Self-consistency (\cref{ch:scaling}) & $k$ samples of one problem & no & whole prompt\\
GRPO rollouts (\cref{ch:grpo}) & $G$ rollouts of one problem & no & whole prompt\\
Distillation training (\cref{ch:distill}) & different (prompt, trace) pairs & yes (right) & no\\
\bottomrule
\end{tabularx}
\end{table}

\subsection{Batched benchmark evaluation}
Evaluating on MATH-500 means generating answers to 500 unrelated problems. With a batch size of
$B$, the evaluator tokenizes $B$ prompts, left-pads them, generates until every row has emitted EOS
or hit the token limit, trims each row at its EOS, and hands the texts to the verifier. The only new knob is the batch size:
\begin{lstlisting}
for i in range(0, len(problems), batch_size):
    chunk = problems[i:i + batch_size]
    prompts = [tok.encode(render_prompt(p["problem"])) for p in chunk]
    outs = generate_batch(model, prompts, max_new_tokens=2048,
                          pad_id=tok.pad_token_id, eos_id=tok.pad_token_id)
    results += [grade(tok.decode(o), p["answer"]) for o, p in zip(outs, chunk)]
\end{lstlisting}
The best batch size is hardware-dependent: increase it until either throughput stops improving or
memory runs out (reasoning traces of thousands of tokens make the KV cache, not the weights, the
limit). Because problems are sorted by nothing in particular, a batch's run time is set by its
\emph{longest} answer; sorting problems by prompt length reduces padding, but output lengths are
unknown in advance, which is the deeper reason servers need continuous batching.

\subsection{Batched self-consistency sampling}
Self-consistency (\cref{ch:scaling}) samples $k$ answers to the \emph{same} prompt at nonzero
temperature and takes a majority vote. Here the batch dimension holds the $k$ samples: we repeat the
prompt $k$ times, so every row has identical length and no padding or special mask is needed. The
plain model class works as is. Each row needs its own random stream, which \code{torch.multinomial}
provides naturally when sampling a $(k,V)$ probability matrix.

\begin{systemslens}[Do not prefill the same prompt $k$ times]
Repeating the prompt $k$ times also repeats its prefill $k$ times and stores $k$ identical copies of
its KV cache. The obvious fix is to prefill once and \emph{fork} the cache: all samples share the
prompt's KV blocks and only allocate private blocks for their own continuations. With paged memory
(\cref{sec:serving:paged}) this is a reference-count increment per block, and a copy-on-write if a
sample must write into a partially filled shared block. For a 1,000-token problem and $k=16$
samples of 500 tokens each, sharing cuts the KV footprint from $16\times1500=24{,}000$ to
$1000+16\times500=9{,}000$ tokens and the prefill work by $16\times$. Lab~22 measures this.
\end{systemslens}

\subsection{Batched GRPO rollouts}
GRPO (\cref{ch:grpo}) samples a \emph{group} of $G$ rollouts per prompt and compares their rewards.
Again the rows share one prompt, so batching rollouts needs no padding logic, and the same
prompt-sharing argument applies. The batch size for rollouts trades throughput against memory:
generating 8 rollouts of 512 tokens in parallel needs 8 KV caches; if memory is tight, generate the
group in several smaller batches. Self-refinement (\cref{ch:refine}) is the counter-example in Part~VI:
each round depends on the previous one, so it is inherently sequential per problem and only benefits
from batching \emph{across} problems, which requires more bookkeeping.

\subsection{Batched distillation training}
Distillation (\cref{ch:distill}) fine-tunes a student on (prompt, teacher reasoning trace) pairs of
widely varying length, so we are back to batching unrelated sequences, this time for \emph{training}.
Sequences are right-padded, pad positions get loss label $-100$, and the attention mask hides pad keys.
Training also stores activations for the backward pass, so memory per row is far higher than in
evaluation, and the useful batch size is smaller. When lengths vary a lot, \emph{packing} several
short examples into one row with block-diagonal attention (\cref{ch:instruct}) beats padding.

% =====================================================================
\section{Single-sequence versus batched execution}\label{sec:serving:tradeoffs}

Batched execution always costs more memory; whether it saves time depends on the shape of the work.
Two quantities predict the outcome. The first is \emph{padding efficiency}: for a batch with prompt
lengths $\ell_1,\ldots,\ell_B$ padded to $\ell_{\max}$,
\begin{equation}
  \eta_{\text{pad}} = \frac{\sum_b \ell_b}{B\,\ell_{\max}},
\end{equation}
and the analogous quantity for output lengths when the batch runs until its longest row finishes.
The second is how much of the step time is amortizable, which the roofline tells us.

\begin{napkinmath}[When batching barely helps]
Take 64 MATH problems with prompt lengths uniform in $[60,240]$ and output lengths uniform in
$[50,1500]$. Expected padding efficiency of the prompts is $\approx150/240=0.63$; for outputs, a batch of
64 runs to $\approx\max\approx1480$ steps while the average row needs 775, so $\eta\approx0.52$: half of
all decode slots are wasted on finished rows. On a GPU the weight-read amortization ($\approx64\times$)
still wins handily. On a CPU running a small model, where step time is roughly proportional to rows,
the same batch can be \emph{slower} than sequential execution. And in a training job that is already
compute-bound, batching adds padding work without saving any memory traffic.
\end{napkinmath}

Our own CPU measurement in \cref{fig:inference:measured} showed sub-linear scaling (about $9\times$
throughput at $B=32$), and published GPU comparisons show large wall-clock reductions from batching
evaluation and rollout generation, with smaller gains for distillation training, where the GPU is
already busy with large matmuls at batch size one. The general rule: batch when the hardware is
underutilized, and expect the gain to shrink as per-row work (padding, KV reads, activations) grows.

% =====================================================================
\section{From static to continuous batching}\label{sec:serving:continuous}

Everything so far was \emph{offline}: we knew all prompts in advance. A server faces an \emph{online}
problem: requests arrive at random times with unknown output lengths, and each user cares about their
own latency. The batching we have used so far, called \emph{static} or request-level batching, waits
until $B$ requests are available (or a timeout fires), runs them as a rectangle until the last one
finishes, then starts the next batch. Three things go wrong:
\begin{enumerate}
  \item \textbf{Head-of-line blocking.} A request arriving just after a batch started waits for the
    entire batch, including its longest member, before its own prefill begins: TTFT explodes.
  \item \textbf{Idle slots.} Rows that finish early sit idle until the batch ends
    (\cref{fig:serving:eos}, strategy 2).
  \item \textbf{Padding and reservation.} Prompts are padded to the longest, and KV memory must be
    reserved for the maximum possible output length.
\end{enumerate}

\textbf{Iteration-level scheduling}, introduced by Orca (Yu et al., 2022) and now called
\emph{continuous batching}, makes scheduling decisions at every forward pass instead of every batch.
After each step, finished sequences leave and waiting requests join; the batch is a set of
sequences, each at its own position (\cref{fig:serving:timeline}).

\begin{figure}[ht]
\centering
\begin{tikzpicture}[x=0.42cm,y=0.5cm,
  pf/.style={fill=fcompute!55,draw=white},
  dc/.style={fill=napkin!45,draw=white},
  wt/.style={fill=black!8,draw=white},
  idle/.style={pattern=north east lines,pattern color=black!35,draw=white}]
\node[font=\small\bfseries,anchor=west] at (0,5.3) {static batching ($B=3$)};
% requests: R1 len 4, R2 len 9, R3 len 6 start at 0 ; R4 arrives t=2 (len 3), R5 arrives t=5 (len 5)
\foreach \y/\l in {4/R1,3/R2,2/R3,1/R4,0/R5} \node[font=\scriptsize,anchor=east] at (-0.2,\y+0.5) {\l};
% batch 1: prefill t0-2, decode to 2+9=11
\fill[pf] (0,4) rectangle (2,5); \fill[dc] (2,4) rectangle (6,5); \fill[idle] (6,4) rectangle (11,5);
\fill[pf] (0,3) rectangle (2,4); \fill[dc] (2,3) rectangle (11,4);
\fill[pf] (0,2) rectangle (2,3); \fill[dc] (2,2) rectangle (8,3); \fill[idle] (8,2) rectangle (11,3);
% batch 2 starts at 11
\fill[wt] (2,1) rectangle (11,2); \fill[pf] (11,1) rectangle (13,2); \fill[dc] (13,1) rectangle (16,2); \fill[idle] (16,1) rectangle (18,2);
\fill[wt] (5,0) rectangle (11,1); \fill[pf] (11,0) rectangle (13,1); \fill[dc] (13,0) rectangle (18,1);
\draw[-{Stealth}] (0,-0.4) -- (19.5,-0.4) node[right,font=\scriptsize]{steps};
\draw[warnred,thick] (2,1) -- (2,2.15); \draw[warnred,thick] (5,0) -- (5,1.15);
\begin{scope}[xshift=0cm,yshift=-3.9cm]
\node[font=\small\bfseries,anchor=west] at (0,5.3) {continuous batching (max 3 running)};
\foreach \y/\l in {4/R1,3/R2,2/R3,1/R4,0/R5} \node[font=\scriptsize,anchor=east] at (-0.2,\y+0.5) {\l};
\fill[pf] (0,4) rectangle (2,5); \fill[dc] (2,4) rectangle (6,5);
\fill[pf] (0,3) rectangle (2,4); \fill[dc] (2,3) rectangle (11,4);
\fill[pf] (0,2) rectangle (2,3); \fill[dc] (2,2) rectangle (8,3);
\fill[wt] (2,1) rectangle (6,2); \fill[pf] (6,1) rectangle (7,2); \fill[dc] (7,1) rectangle (10,2);
\fill[wt] (5,0) rectangle (8,1); \fill[pf] (8,0) rectangle (9,1); \fill[dc] (9,0) rectangle (14,1);
\draw[-{Stealth}] (0,-0.4) -- (19.5,-0.4) node[right,font=\scriptsize]{steps};
\draw[warnred,thick] (2,1) -- (2,2.15); \draw[warnred,thick] (5,0) -- (5,1.15);
\end{scope}
\begin{scope}[yshift=-5.0cm]
\fill[pf] (0,0) rectangle (1,0.7); \node[font=\scriptsize,anchor=west] at (1.2,0.35) {prefill};
\fill[dc] (4,0) rectangle (5,0.7); \node[font=\scriptsize,anchor=west] at (5.2,0.35) {decode};
\fill[wt] (8,0) rectangle (9,0.7); \node[font=\scriptsize,anchor=west] at (9.2,0.35) {queued};
\fill[idle] (12,0) rectangle (13,0.7); \node[font=\scriptsize,anchor=west] at (13.2,0.35) {finished, still in batch};
\draw[warnred,thick] (21.5,-0.05) -- (21.5,0.75); \node[font=\scriptsize,anchor=west] at (21.8,0.35) {arrival};
\end{scope}
\end{tikzpicture}
\caption{Five requests under static and continuous batching. Statically, R4 and R5 wait for the
whole first batch (their TTFT includes R2's full generation) and finished rows idle. Continuously,
R4 joins as soon as R1 leaves and R5 as soon as R3 leaves; the whole workload finishes at step 14
instead of 18. (In a real engine the joiners' prefills share steps with others' decodes; see
\cref{sec:serving:chunked}.)}
\label{fig:serving:timeline}
\end{figure}

Continuous batching needs two capabilities that the rectangle formulation lacks. First, the forward
pass must handle a \emph{ragged} batch: some sequences are in prefill with hundreds of new tokens,
others decode one token each, and each attends over its own history. Engines therefore flatten all
scheduled tokens into one $(N_{\text{tok}},d)$ matrix, run every dense layer as one GEMM over all of
them (no padding at all), and only split by sequence inside attention, using variable-length
attention kernels that take per-sequence offsets (``varlen'' or ``cu\_seqlens'' interfaces in
FlashAttention, FlashInfer). Second, KV memory must be allocated and freed at fine granularity as
sequences come and go, which is the subject of the next section.

\begin{systemslens}[Why flattening beats padding]
In a flattened batch the dense layers cost $2N\cdot\sum_s n_s$ FLOPs where $n_s$ is the number of
\emph{new} tokens of sequence $s$ this step, exactly the useful work. Padding to a rectangle would cost
$2N\cdot S\cdot\max_s n_s$: a single 2,000-token prefill in a batch of 63 decodes would inflate the
GEMMs by $64\times$. Flattening also makes the LM head cheap: it is evaluated only on the last new
token of each sequence ($S$ rows, not $\sum_s n_s$).
\end{systemslens}

% =====================================================================
\section{Paged KV memory}\label{sec:serving:paged}

Before 2023, serving systems stored each sequence's KV cache in a contiguous tensor sized for the
maximum sequence length. Kwon et al.\ (2023) measured that such systems used only a small fraction
of their KV memory for actual tokens and proposed \textbf{PagedAttention}, the core of vLLM. The
waste has three sources:
\begin{itemize}
  \item \textbf{Reservation}: memory for future tokens is reserved at admission, although the output
    length is unknown and usually far below the maximum.
  \item \textbf{Internal fragmentation}: unused tail of each reserved region.
  \item \textbf{External fragmentation}: free memory split into gaps too small for a new contiguous
    allocation.
\end{itemize}

\begin{napkinmath}[How much memory reservation wastes]
Reserve 2,048 tokens per request; actual lengths average 300. Utilization is $300/2048\approx15\%$:
85\% of the KV pool is dead weight, and batch size is $6.8\times$ smaller than it could be. With paging
in 16-token blocks, only the last block of each sequence is partially filled, wasting on average 7.5
tokens per sequence: $7.5/307.5\approx2.4\%$. Lab~19 (Task~4) computes both for any length distribution.
\end{napkinmath}

The solution is borrowed from operating-system virtual memory (\cref{fig:serving:paged}). KV memory
is a pool of fixed-size \emph{blocks} (pages), each holding the K and V vectors of \texttt{block\_size}
tokens for all layers and KV heads. Each sequence has a \emph{block table} mapping logical block $i$
(tokens $i\cdot\text{bs}$ to $(i+1)\cdot\text{bs}-1$) to a physical block id. Blocks are allocated
on demand when a sequence crosses a block boundary and returned to a free list when it finishes. The
attention kernel follows the table: the key of logical position $p$ lives at physical slot
\begin{equation}\label{eq:serving:slot}
  \text{slot}(p) = \text{table}[\lfloor p/\text{bs}\rfloor]\cdot\text{bs} + (p \bmod \text{bs}).
\end{equation}

\begin{figure}[ht]
\centering
\begin{tikzpicture}[
  tk/.style={draw=black!35,minimum width=0.52cm,minimum height=0.45cm,font=\tiny,inner sep=0pt},
  pb/.style={draw=black!50,minimum width=2.9cm,minimum height=0.45cm,font=\tiny,inner sep=0pt},
  arr/.style={-{Stealth[length=1.4mm]},thick}]
% ---- sequence A: logical blocks L0 L1 L2
\node[font=\scriptsize\bfseries,anchor=west] at (-0.35,0.75) {seq A (10 tokens)};
\foreach \g/\col/\toks in {0/fcompute!25/{The,sky,is,blue},1/fdata!25/{be,cau,se,of},2/napkin!35/{Ray,leigh,{},{}}} {
  \foreach \t [count=\k from 0] in \toks {
    \node[tk,fill=\col] at (2.3*\g+0.52*\k,0) {\t};
  }
  \node[font=\tiny,text=black!60] at (2.3*\g+0.78,-0.42) {logical block \g};
}
\node[font=\scriptsize,anchor=south] at (8.0,0.3) {block table A};
\foreach \v/\col [count=\k from 0] in {5/fcompute!25,1/fdata!25,6/napkin!35}
  \node[tk,fill=\col,font=\scriptsize] (A\k) at (7.48+0.52*\k,0) {\v};
% ---- sequence B
\node[font=\scriptsize\bfseries,anchor=west] at (-0.35,-1.25) {seq B (6 tokens)};
\foreach \g/\col/\toks in {0/labpurple!25/{Two,plus,two,is},1/keygreen!30/{four,.,{},{}}} {
  \foreach \t [count=\k from 0] in \toks {
    \node[tk,fill=\col] at (2.3*\g+0.52*\k,-2.0) {\t};
  }
  \node[font=\tiny,text=black!60] at (2.3*\g+0.78,-2.42) {logical block \g};
}
\node[font=\scriptsize,anchor=north] at (7.74,-2.3) {block table B};
\foreach \v/\col [count=\k from 0] in {2/labpurple!25,7/keygreen!30}
  \node[tk,fill=\col,font=\scriptsize] (B\k) at (7.48+0.52*\k,-2.0) {\v};
% ---- physical pool
\node[font=\scriptsize\bfseries] at (12.0,1.15) {physical KV pool};
\foreach \p/\col/\lab in {0/white/free,1/fdata!25/{A: be cau se of},2/labpurple!25/{B: Two plus two is},3/white/free,4/white/free,5/fcompute!25/{A: The sky is blue},6/napkin!35/{A: Ray leigh -- --},7/keygreen!30/{B: four . -- --}} {
  \node[pb,fill=\col] (P\p) at (12.0,0.6-0.52*\p) {\lab};
  \node[font=\tiny,anchor=east] at (10.45,0.6-0.52*\p) {\p};
}
\draw[arr,fcompute] (A0.south) to[out=-70,in=180] (P5.west);
\draw[arr,fdata!80!black] (A1.north) to[out=70,in=180] (P1.west);
\draw[arr,napkin] (A2.south) to[out=-50,in=180] (P6.west);
\draw[arr,labpurple] (B0.north) to[out=80,in=180] (P2.west);
\draw[arr,keygreen!80!black] (B1.south) to[out=-40,in=180] (P7.west);
\end{tikzpicture}
\caption{Paged KV memory. Logically contiguous sequences are stored in non-contiguous physical
blocks; a per-sequence block table translates positions to slots (\cref{eq:serving:slot}). Only the
last block of each sequence can be partially empty, and any free block can serve any sequence.}
\label{fig:serving:paged}
\end{figure}

\paragraph{Choosing the block size.} Small blocks waste less memory and allow finer sharing; large
blocks make the attention kernel's gathers more efficient (longer contiguous reads, fewer table
lookups) and shrink block tables. Engines typically use 16 tokens per block on GPUs. For Qwen3-8B a
16-token block holds $16\times\SI{144}{KiB}=\SI{2.25}{MiB}$ across all layers; a \SI{20}{GB} pool holds about
8,500 blocks, i.e.\ 136k tokens.

\begin{systemslens}[The cost of indirection]
Paging turns the attention kernel's KV reads from one contiguous stream into a gather over blocks.
On a GPU this is nearly free if each block is a large contiguous chunk (each warp reads whole blocks)
and the block table is small enough to sit in shared memory or registers. The PagedAttention paper
reports only modest kernel-level overhead in exchange for several-fold higher batch sizes. Later
systems (e.g.\ vAttention) instead use CUDA virtual-memory APIs to keep the KV cache virtually
contiguous while mapping physical pages on demand, so unmodified attention kernels can be used.
\end{systemslens}

% =====================================================================
\section{Sharing KV state: prefix caching and copy-on-write}\label{sec:serving:prefix}

Once KV memory is managed in blocks with reference counts, \emph{sharing} becomes cheap: two
sequences whose block tables point to the same physical block simply increment its count. Three
patterns are ubiquitous:
\begin{itemize}
  \item \textbf{Parallel sampling} ($k$ samples or $G$ rollouts of one prompt): fork the prompt's
    blocks after one prefill.
  \item \textbf{Shared system prompts and few-shot templates}: thousands of requests start with the
    same hundreds of tokens.
  \item \textbf{Multi-turn chat and agents}: turn $t+1$ resends the full conversation of turn $t$ plus a
    few new tokens. Without caching, the whole history is prefilled again every turn.
\end{itemize}

\paragraph{Hash-based automatic prefix caching.} vLLM identifies every \emph{full} block by a hash of
its tokens \emph{and} the hash of the previous block, so a hash identifies the entire prefix up to and
including that block. Completed blocks are published in a hash table. A new request walks its prompt
block by block, computing chained hashes, and reuses every block it finds until the first miss; it
then only prefills the remainder. Blocks whose reference count drops to zero are not freed; they stay
in the table as \emph{evictable} and are reclaimed in least-recently-used order when the allocator runs
dry. One subtlety: at least one prompt token must still be computed, because the model has to produce
the logits for the first output token.

\begin{lstlisting}[caption={Prefix lookup with chained block hashes (llmsys.serving.block\_manager).},label={lst:serving:lookup}]
def block_hash(prev_hash, tokens):
    return hash((prev_hash, tuple(tokens)))      # identifies the whole prefix

def lookup_prefix(self, tokens):
    bs, out, h = self.block_size, [], None
    for i in range((len(tokens) - 1) // bs):     # leave >= 1 token to compute
        h = block_hash(h, tokens[i * bs:(i + 1) * bs])
        b = self.hash_to_block.get(h)
        if b is None:
            break                                # first miss ends the match
        if self.ref_count[b] == 0:
            self.evictable.pop(b, None)          # revive a cached block
        self.ref_count[b] += 1
        out.append(b)
    return out
\end{lstlisting}

\paragraph{RadixAttention.} SGLang (Zheng et al., 2024) organizes cached prefixes as a \emph{radix
tree} keyed by token sequences (\cref{fig:serving:radix}). Each edge holds a run of tokens and the
corresponding KV pages; a request matches the longest path from the root; insertion splits edges where
sequences diverge; LRU eviction removes leaves first. The tree makes the cache structure explicit,
supports matching at token rather than block granularity, and enables \emph{cache-aware scheduling}:
ordering the waiting queue so requests that share long prefixes run close together, which raises the hit
rate. Hash chains and radix trees are two indexes over the same idea.

\begin{figure}[ht]
\centering
\begin{tikzpicture}[
  nd/.style={draw,rounded corners=2pt,font=\scriptsize,align=center,inner sep=3pt},
  lv/.style={nd,fill=napkin!20},
  arr/.style={-{Stealth[length=1.4mm]},thick,black!60},
  level distance=1.35cm]
\node[nd,fill=black!8] (root) at (0,0) {root};
\node[nd,fill=fcompute!15] (sys) at (0,-1.3) {system prompt (512 tok)\\ref = 4};
\node[nd,fill=fdata!15] (fs) at (-3.6,-2.7) {few-shot math\\examples (900)};
\node[nd,fill=fdata!15] (chat) at (3.4,-2.7) {``You are a helpful\\assistant'' (40)};
\node[lv] (q1) at (-5.3,-4.2) {Q: solve\ldots\\(60) + answer};
\node[lv] (q2) at (-2.0,-4.2) {Q: prove\ldots\\(85) + answer};
\node[nd,fill=fdata!15] (t1) at (2.2,-4.2) {turn 1\\(user + assistant)};
\node[lv] (u2) at (4.8,-4.2) {other user\\turn 1};
\node[lv] (t2) at (2.2,-5.6) {turn 2 (new tokens only\\need prefill)};
\foreach \a/\b in {root/sys,sys/fs,sys/chat,fs/q1,fs/q2,chat/t1,chat/u2,t1/t2} \draw[arr] (\a) -- (\b);
\node[font=\scriptsize,align=left,anchor=west,text=black!70] at (-6.3,-0.5)
  {match: longest root-to-node path\\insert: split edge where tokens diverge\\evict: LRU leaves with ref = 0};
\end{tikzpicture}
\caption{A radix tree of cached prefixes. Requests sharing the system prompt reuse its KV pages;
the follow-up turn of a conversation matches everything up to and including turn~1, so only the new
user message is prefilled. Leaves (orange) are the first candidates for eviction.}
\label{fig:serving:radix}
\end{figure}

\paragraph{Copy-on-write.} Sharing full blocks is safe because full blocks are immutable: new tokens are
always appended to the \emph{last} block. When a forked sequence (parallel sampling, beam search) shares a
\emph{partially filled} last block, the first sequence to append must copy it into a private block
first, exactly like copy-on-write pages after \code{fork()} in an operating system.

\begin{napkinmath}[What prefix caching buys]
In our simulator (\cref{sec:serving:sim}), 600 requests to Qwen3-8B on an H100 at 8 requests/s, each with a
1,024-token system prompt drawn from four variants, plus a unique 64--512-token question: prefix caching
reused 610k of 787k prompt tokens (78\%), cut the prefill work by $4.5\times$, and reduced mean TTFT from
\SI{52}{ms} to \SI{19}{ms}. Throughput was unchanged because the system was not saturated at that load;
under saturation, saved prefill compute turns directly into throughput.
\end{napkinmath}

% =====================================================================
\section{Chunked prefill and the token budget}\label{sec:serving:chunked}

Continuous batching mixes prefills and decodes in the same step, which creates a new problem. A
4,000-token prefill takes tens of milliseconds of compute; every decoding user sharing that step waits
for it, so their inter-token latency spikes (a ``generation stall''). Not admitting prefills while
decodes run would starve new requests instead. Sarathi-Serve (Agrawal et al., 2024) resolves this with
\textbf{chunked prefill} under a per-step \textbf{token budget} (\cref{fig:serving:chunked}):
\begin{enumerate}
  \item Each step processes at most $\tau$ tokens in total (e.g.\ 512 or 2,048).
  \item Running decodes are scheduled first, one token each.
  \item The remaining budget is filled with \emph{chunks} of pending prefills; a long prompt is split
    across several steps, carrying its KV cache forward between chunks.
\end{enumerate}
Every step then has roughly the same, bounded cost: decode latency becomes predictable, and the
prefill chunks ``piggyback'' on decode steps that would otherwise underuse the GPU's compute. The
budget $\tau$ is a dial between TTFT (large $\tau$ finishes prefills sooner) and TPOT (small $\tau$ keeps
steps short); it is usually chosen near the point where a step becomes compute-bound.

\begin{figure}[ht]
\centering
\begin{tikzpicture}[x=0.62cm,y=0.022cm]
\foreach \off/\title in {0/{no chunking},11/{chunked prefill, budget 64}} {
  \begin{scope}[xshift=\off*0.62cm]
  \draw[-{Stealth}] (0,0) -- (9.3,0) node[below left,font=\scriptsize]{step};
  \draw[-{Stealth}] (0,0) -- (0,205);
  \node[font=\scriptsize\bfseries,anchor=south] at (4.5,212) {\title};
  \draw[dashed,warnred] (0,64) -- (9,64);
  \foreach \yy in {0,64,128,192} \node[font=\tiny,anchor=east] at (-0.1,\yy) {\yy};
  \end{scope}
}
\node[font=\scriptsize,rotate=90] at (-1.3,100) {tokens processed in step};
\node[font=\scriptsize,text=warnred,anchor=west] at (20.4,64) {$\tau$};
% no chunking: decodes 6 per step, step 3 has 180 prefill
\foreach \s/\d/\p in {1/6/0,2/6/0,3/6/180,4/7/0,5/7/0,6/7/0,7/7/0,8/7/0} {
  \fill[napkin!60] (\s-0.35,0) rectangle (\s+0.35,\d);
  \fill[fcompute!60] (\s-0.35,\d) rectangle (\s+0.35,\d+\p);
}
\node[font=\tiny,align=center] at (5.6,175) {one slow step:\\every decoding\\user waits};
% chunked
\foreach \s/\d/\p in {1/6/0,2/6/0,3/6/58,4/6/58,5/6/58,6/7/6,7/7/0,8/7/0} {
  \fill[napkin!60] (11+\s-0.35,0) rectangle (11+\s+0.35,\d);
  \fill[fcompute!60] (11+\s-0.35,\d) rectangle (11+\s+0.35,\d+\p);
}
\node[font=\tiny,align=center] at (16.5,120) {58 + 58 + 58 + 6\\prefill tokens ride\\along with decodes};
\fill[napkin!60] (2,-40) rectangle (2.7,-25); \node[font=\scriptsize,anchor=west] at (2.8,-32) {decode tokens (one per running sequence)};
\fill[fcompute!60] (13,-40) rectangle (13.7,-25); \node[font=\scriptsize,anchor=west] at (13.8,-32) {prefill tokens};
\end{tikzpicture}
\caption{A 180-token prompt arrives while 6--7 sequences decode. Without chunking, step 3 processes 186
tokens and every decoding user sees one very slow token. With a 64-token budget the prompt is split into
chunks of 58, 58, 58 and 6 tokens that ride along with the decodes (the prompt's first token appears
one step later, but no step exceeds the budget).}
\label{fig:serving:chunked}
\end{figure}

% =====================================================================
\section{Scheduling, admission, and preemption}\label{sec:serving:scheduling}

The scheduler decides, every step, which sequences run and how many tokens each processes. Its inputs
are the waiting queue, the running set, the free block count, and the token budget. A minimal policy,
the one our engine implements, is:
\begin{enumerate}
  \item \textbf{Serve running sequences first} (oldest first), giving each $\min(\text{remaining},
    \text{budget})$ tokens: 1 for decodes, a chunk for unfinished prefills.
  \item If a running sequence needs a new block and none is free, \textbf{preempt} the \emph{youngest}
    running sequence (last admitted, least work lost) until space is available.
  \item With the remaining budget, \textbf{admit} waiting requests first-come-first-served, after a
    prefix-cache lookup, as long as their first chunk fits in memory and the running-set limit allows.
    Do not admit in a step that had to preempt (memory is tight).
\end{enumerate}

\paragraph{Preemption: swap or recompute?} A preempted sequence loses its KV blocks. \emph{Swapping}
copies them to host memory over PCIe and back later; \emph{recomputation} discards them and later
re-prefills the prompt \emph{plus the tokens generated so far}. Per token, swapping costs
$2\kappa/\beta_{\text{PCIe}}$ (out and back in) while recomputation costs about $2N/\Pi$ plus attention.
For Qwen3-8B ($\kappa=\SI{144}{KiB}$) over \SI{50}{GB/s} that is \SI{5.9}{\micro s} versus
\SI{16.6}{\micro s} at peak FLOP/s: with GQA's small caches, swapping is competitive. For a model with
full multi-head attention and $4\times$ larger caches, recomputation wins up to tens of thousands of
tokens. In practice, swapping many small blocks over PCIe is inefficient, swap traffic competes with
other host transfers, and recomputation needs no host memory and may hit the prefix cache, so engines
often default to recomputation. Our engine recomputes; Lab~19's tests verify that preempted sequences produce exactly the
same tokens as unpreempted ones.

\begin{systemslens}[Beyond FCFS]
FCFS is fair in arrival order but ignores everything else. Production schedulers add: priorities and
SLO classes (interactive vs.\ batch); \emph{shortest-job-first} approximations using predicted output
length to cut mean latency; fairness across tenants measured in served tokens (Virtual Token Counter,
Sheng et al., 2024) to stop one heavy user from starving others; cache-aware ordering to maximize
prefix hits; and admission control that rejects or queues requests early when the predicted TTFT would
violate the SLO, which protects goodput rather than raw throughput.
\end{systemslens}

% =====================================================================
\section{Disaggregating prefill and decode}\label{sec:serving:disagg}

Chunked prefill makes prefill and decode coexist; disaggregation separates them. DistServe (Zhong et
al., 2024) and Splitwise (Patel et al., 2024) run prefill and decode on \emph{different} GPU pools: a
prefill instance computes the prompt's KV cache and ships it to a decode instance, which then streams
tokens (\cref{fig:serving:disagg}). The motivation is the asymmetry of \cref{sec:inference:roofline}:
\begin{itemize}
  \item Prefill is compute-bound and its latency (TTFT) benefits from tensor parallelism and small
    batches; decode is bandwidth-bound and benefits from huge batches.
  \item Co-locating them couples their SLOs: any prefill in a step delays every decode in it.
  \item Separate pools can be sized, parallelized and even built from different hardware
    independently, maximizing goodput per GPU.
\end{itemize}

\begin{figure}[ht]
\centering
\begin{tikzpicture}[
  gp/.style={draw,rounded corners=3pt,minimum width=3.0cm,minimum height=1.6cm,align=center,font=\scriptsize},
  arr/.style={-{Stealth[length=1.8mm]},thick}]
\node[draw,rounded corners=2pt,fill=black!5,font=\scriptsize,align=center,minimum height=1.0cm] (rt) at (0,0) {router /\\scheduler};
\node[gp,fill=fcompute!15] (pf) at (4.0,0.95) {\textbf{prefill pool}\\compute-bound\\TP for low TTFT};
\node[gp,fill=napkin!20] (dc) at (4.0,-1.35) {\textbf{decode pool}\\bandwidth-bound\\large batches};
\node[draw,rounded corners=2pt,fill=fmemory!15,font=\scriptsize,align=center] (kv) at (8.6,0.95) {KV store\\(HBM / DRAM / SSD)};
\draw[arr] (-2.0,0) node[left,font=\scriptsize]{requests} -- (rt);
\draw[arr] (rt) -- node[above,font=\tiny,sloped]{prompt} (pf.west);
\draw[arr,fcomm] (pf.south) -- node[right,font=\scriptsize,text=fcomm,align=left]{KV transfer\\(NVLink / RDMA)} (dc.north);
\draw[arr] (dc.west) -- node[below,font=\tiny,sloped]{tokens} (rt.south east);
\draw[arr,dashed,black!50] (pf.east) -- node[above,font=\tiny]{publish} (kv.west);
\draw[arr,dashed,black!50] (kv.south) |- node[pos=0.25,right,font=\tiny]{reuse} (dc.east);
\end{tikzpicture}
\caption{Disaggregated serving. Prefill and decode run on separately sized pools; the KV cache moves
between them over a fast interconnect. Systems such as Mooncake add a distributed KV store so that
caches computed once can be reused across requests and machines.}
\label{fig:serving:disagg}
\end{figure}

\begin{napkinmath}[Is the KV transfer affordable?]
A 2,000-token prompt on Qwen3-8B produces $2000\times\SI{144}{KiB}\approx\SI{295}{MB}$ of KV cache. Its prefill takes
at least $2000\cdot\SI{15.1}{GFLOP}/\SI{989}{TFLOP/s}\approx\SI{31}{ms}$ on an H100. Shipping
\SI{295}{MB} takes \SI{0.33}{ms} over NVLink (\SI{900}{GB/s}) and $\approx\SI{5.9}{ms}$ over a 400\,Gb/s
RDMA NIC (\SI{50}{GB/s}): well below the prefill time, and it can be pipelined layer by layer as
the prefill progresses. The transfer only becomes the bottleneck for models with very large per-token
caches on slow networks, which is one more reason for GQA/MLA (\cref{ch:modern}) and fp8 KV caches.
\end{napkinmath}

\paragraph{Tiered KV caches.} Prefix caches outgrow HBM quickly. Engines therefore spill evicted blocks
to host DRAM or SSD and reload them on a hit, provided reloading is faster than recomputing:
at \SI{25}{GB/s} PCIe, reloading the \SI{295}{MB} prefix above takes \SI{12}{ms}, cheaper than its
\SI{31}{ms} prefill. Mooncake (Qin et al., 2024) builds a cluster-wide KV store around this trade-off.

% =====================================================================
\section{Anatomy of a production engine}\label{sec:serving:anatomy}

Before building our own engine, it helps to see what a production engine adds around the scheduler
and the KV cache. Three pieces deserve a systems person's attention.

\paragraph{Sizing the KV pool at start-up.} Engines do not ask the user how many KV blocks to
allocate. At start-up they load the weights, run a \emph{profiling} forward pass with the largest
batch the scheduler may ever issue (e.g.\ \code{max\_batched\_tokens} tokens of prefill), record the
peak activation memory, and give everything that remains, up to a utilization fraction, to the KV pool.

\begin{napkinmath}[KV pool for Qwen3-8B on one H100]
Usable HBM at 90\% utilization: $0.9\times\SI{80}{GB}=\SI{72}{GB}$. Weights: \SI{16.4}{GB}. Peak activations of an
8,192-token profiling step: the largest transient is the MLP intermediate, $8192\times12288\times2$ bytes
$\approx\SI{0.2}{GB}$ per tensor; a few such tensors plus logits for the sampled rows, the CUDA context
and CUDA-graph memory pools come to perhaps \SI{4}{GB}. KV pool:
$72-16.4-4\approx\SI{51.6}{GB}$, i.e.\ $51.6\,\text{GB}/(16\times\SI{144}{KiB})\approx21{,}900$ blocks
$\approx\num{350000}$ tokens. At an average of 2,000 tokens per sequence, that is room for $\approx175$
concurrent sequences, comfortably past the batch size where decode throughput saturates in
\cref{tab:inference:decode}. For Qwen3-32B on the same GPU, weights alone take \SI{65.5}{GB} and the KV
pool shrinks to a few GB: such a model is served with tensor parallelism across GPUs largely to gain
KV memory, not FLOPs (\cref{ch:efficient}).
\end{napkinmath}

\paragraph{Process architecture.} Python is single-threaded, and the GPU must never wait for it. Modern
engines therefore split work across processes (\cref{fig:serving:procs}): an API server handles HTTP,
tokenization, chat templates and detokenization; an \emph{engine core} process runs the scheduler and
the step loop; GPU worker processes (one per GPU for tensor parallelism) execute the model. The
processes communicate over shared memory or ZeroMQ sockets, so CPU-heavy work (tokenizing a 100k-token
prompt, detokenizing thousands of streams, JSON encoding) overlaps with GPU execution. Inside the
engine core, \emph{asynchronous scheduling} prepares the inputs of step $t+1$ (scheduling, block
allocation, building the slot mapping on the host and copying it to the device) while the GPU still
executes step $t$, so the GPU sees back-to-back steps.

\begin{figure}[ht]
\centering
\begin{tikzpicture}[
  pr/.style={draw,rounded corners=3pt,minimum height=1.25cm,align=center,font=\scriptsize,minimum width=2.9cm},
  arr/.style={{Stealth[length=1.5mm]}-{Stealth[length=1.5mm]},thick,black!60}]
\node[font=\scriptsize] (cl) at (-1.6,0) {clients};
\node[pr,fill=fdata!12] (api) at (1.6,0) {\textbf{API server}\\HTTP / SSE, chat template,\\tokenize, detokenize};
\node[pr,fill=sysblue!10] (core) at (5.6,0) {\textbf{engine core}\\scheduler, block manager,\\async step loop};
\node[pr,fill=fcompute!15,minimum width=2.3cm] (w0) at (9.6,0.8) {\textbf{GPU worker 0}\\TP rank 0};
\node[pr,fill=fcompute!15,minimum width=2.3cm] (w1) at (9.6,-0.8) {\textbf{GPU worker 1}\\TP rank 1};
\draw[arr] (cl) -- (api);
\draw[arr] (api) -- node[above,font=\tiny]{IPC (ZMQ)} (core);
\draw[arr] (core.east) -- (w0.west); \draw[arr] (core.east) -- (w1.west);
\draw[thick,fcomm] (w0.south) -- node[right,font=\tiny,text=fcomm]{NCCL all-reduce} (w1.north);
\node[font=\tiny,text=black!60,align=center] at (3.6,-1.1) {token ids in,\\token ids out};
\node[font=\tiny,text=black!60,align=center] at (7.4,-1.55) {per step: slot mapping,\\block tables, positions};
\end{tikzpicture}
\caption{Process layout of a typical modern engine. CPU-bound request handling, scheduling and GPU
execution run concurrently in separate processes.}
\label{fig:serving:procs}
\end{figure}

\paragraph{Batched sampling with per-request parameters.} Every request may carry its own temperature,
top-$p$, top-$k$, seed, repetition penalty, stop strings and logprob requests. A Python loop over rows (as in our
engine) is fine for a toy but serializes thousands of tiny kernels. Production samplers vectorize: parameters
become per-row tensors, and one sort, one cumulative sum and one \code{multinomial} serve the whole batch
(\cref{lst:serving:batchsample}). Greedy rows are temperature-0 rows handled by a final \code{where}.

\begin{lstlisting}[caption={Vectorized sampling with per-row temperature and top-p.},label={lst:serving:batchsample}]
def sample_batch(logits, temperature, top_p, generator=None):
    # logits: (B, V); temperature, top_p: (B,) tensors. Returns (B,) ids.
    greedy = logits.argmax(dim=-1)
    t = temperature.clamp_min(1e-5)[:, None]
    probs = torch.softmax(logits.float() / t, dim=-1)
    sorted_p, idx = probs.sort(dim=-1, descending=True)
    cum = sorted_p.cumsum(dim=-1)
    sorted_p[(cum - sorted_p) >= top_p[:, None]] = 0.0        # per-row nucleus
    choice = torch.multinomial(sorted_p, 1, generator=generator)  # renormalizes
    sampled = idx.gather(-1, choice).squeeze(-1)
    return torch.where(temperature == 0, greedy, sampled)
\end{lstlisting}

The full sort over $V=151{,}936$ entries per row is the expensive part. Faster samplers avoid it by
rejection sampling or by first pruning to a top-$k$ candidate set, and some draw samples via the
``exponential race'' trick, $\argmax_i\, p_i/E_i$ with $E_i\sim\text{Exp}(1)$, which needs neither a sort
nor a cumulative sum (it samples from $p$ exactly, so top-$p$ must still be applied first). These tricks
matter because at large batch sizes sampling can take a visible share of each decode step.

% =====================================================================
\section{Building a mini engine}\label{sec:serving:engine}

We now assemble the pieces into \ident{llmsys.serving}, a teaching-sized engine of a few hundred lines
that runs our GPT-2 and Qwen3 implementations unchanged (\cref{fig:serving:arch}). It is deliberately
simple (pure PyTorch, one process, a Python loop for per-sequence attention), but its structure mirrors
vLLM's: a scheduler, a block manager, a model runner, and an engine loop.

\begin{figure}[ht]
\centering
\begin{tikzpicture}[
  cmp/.style={draw,rounded corners=2pt,minimum width=3.2cm,minimum height=0.95cm,align=center,font=\scriptsize},
  arr/.style={-{Stealth[length=1.6mm]},thick}]
\node[cmp,fill=black!5] (eng) at (0,0) {\textbf{LLMEngine}\\\code{add\_request / step}};
\node[cmp,fill=sysblue!10] (sch) at (-3.9,-1.7) {\textbf{Scheduler}\\token budget, chunking,\\admission, preemption};
\node[cmp,fill=fmemory!20] (bm) at (-3.9,-3.6) {\textbf{BlockManager}\\free list, ref counts,\\prefix hash table + LRU};
\node[cmp,fill=fcompute!15] (mr) at (3.9,-1.7) {\textbf{ModelRunner}\\ragged forward,\\LM head on last tokens};
\node[cmp,fill=fmemory!20] (pool) at (3.9,-3.6) {\textbf{PagedKVPool}\\$(L,\ \text{slots},H_{kv},d_h)$\\K and V};
\node[cmp,fill=keygreen!12] (seq) at (0,-2.65) {\textbf{Sequence}\\tokens, block table,\\\code{num\_computed}};
\draw[arr] (eng) -- node[above left,font=\tiny]{1. schedule()} (sch);
\draw[arr] (sch) -- node[left,font=\tiny]{alloc / free} (bm);
\draw[arr] (eng) -- node[above right,font=\tiny]{2. forward(chunks)} (mr);
\draw[arr] (mr) -- node[right,font=\tiny]{write, gather} (pool);
\draw[arr,dashed] (eng) -- node[right,font=\tiny,align=left]{3. sample, publish\\blocks, finish} (seq);
\draw[dashed,black!40] (seq) -- (sch); \draw[dashed,black!40] (seq) -- (mr);
\end{tikzpicture}
\caption{Architecture of \ident{llmsys.serving}. The scheduler and block manager deal only in integers
(token counts, block ids); only the model runner touches tensors.}
\label{fig:serving:arch}
\end{figure}

\paragraph{One state variable.} A sequence's tokens are its prompt plus its output so far, and
\ident{num_computed} counts how many of them already have K/V in the pool. Everything follows from
$\text{remaining} = \text{len(tokens)} - \text{num\_computed}$: a fresh request has remaining $=$
prompt length (prefill); a decoding sequence has remaining $=1$ (its last sampled token); a chunked
prefill has remaining $>1$ after its first chunk; a preempted sequence is reset to
$\text{num\_computed}=0$ and simply re-prefills everything; a prefix hit sets $\text{num\_computed}$ to
the matched length. The engine samples a token only when a sequence's chunk reaches the end of its known
tokens. Prefill, decode, chunking, preemption and prefix reuse are thus one code path.

\begin{lstlisting}[caption={The engine step (abridged from llmsys/serving/engine.py).},label={lst:serving:step}]
def step(self):
    batch = self.scheduler.schedule()             # [(seq, n_tokens), ...]
    chunks = [(seq.tokens[seq.num_computed: seq.num_computed + n],
               seq.num_computed, seq.block_table) for seq, n in batch]
    logits = self.runner.forward(chunks)          # one ragged forward pass
    for i, (seq, n) in enumerate(batch):
        seq.num_computed += n
        self._publish_full_blocks(seq)            # feed the prefix cache
        if seq.remaining > 0:
            continue                              # mid-prefill: nothing to sample
        tok = int(sample_next_token(logits[i:i + 1], seq.params.temperature,
                                    seq.params.top_k, seq.params.top_p, seq.generator))
        seq.output.append(tok)
        if seq.is_finished():
            self.scheduler.finish(seq)            # releases (or caches) its blocks
\end{lstlisting}

\paragraph{The ragged forward pass.} \ident{ModelRunner.forward} concatenates the new tokens of all
scheduled sequences into one vector, computes the physical slot of every token with
\cref{eq:serving:slot}, and walks the layers. All projections and the MLP run once on the whole
$(N_{\text{tok}},d)$ matrix. Per layer, the new keys and values are scattered into the pool with
\code{index\_copy\_}; attention then loops over sequences, gathering each one's history through its
block table (\cref{lst:serving:pagedattn}). Finally, the LM head runs only on the last new token of
each sequence. RoPE is applied with per-token positions, so sequences at different positions coexist
in one batch.

\begin{lstlisting}[caption={Reference paged attention for one sequence.},label={lst:serving:pagedattn}]
def paged_attention(q, k_pool, v_pool, block_table, ctx_len, block_size, group_size):
    n_new, H, D = q.shape                       # queries for the last n_new positions
    slots = torch.tensor(slot_mapping(block_table, range(ctx_len), block_size))
    k = k_pool.index_select(0, slots).repeat_interleave(group_size, dim=1)
    v = v_pool.index_select(0, slots).repeat_interleave(group_size, dim=1)
    scores = torch.einsum("qhd,khd->hqk", q.float(), k.float()) / math.sqrt(D)
    past = ctx_len - n_new                      # causal offset for chunked prefill
    mask = torch.triu(torch.ones(n_new, ctx_len, dtype=torch.bool), diagonal=past + 1)
    w = torch.softmax(scores.masked_fill(mask[None], float("-inf")), dim=-1)
    return torch.einsum("hqk,khd->qhd", w, v.float()).to(q.dtype)
\end{lstlisting}

The Python loop over sequences and the explicit gather are what a real engine replaces with a fused
paged-attention kernel: one thread block per (sequence, head), streaming K/V blocks from the table into
shared memory with an online softmax (\cref{ch:attention}), so the gathered K/V never materializes.
Writing that kernel in Triton is a stretch goal in Lab~19 and \cref{app:triton}.

\paragraph{Correctness.} The strongest test of an engine is that it is \emph{invisible}: for greedy
decoding, every request must receive exactly the tokens it would have received alone. Lab~19's tests
run six requests with a shared prefix through the engine under three regimes (plenty of memory; a
5-token budget forcing chunked prefill; a 10-block pool forcing repeated preemption) and compare against
\ident{llmsys.generate.generate} token by token, for both GPT-2 and Qwen3. They also assert the block
manager's invariant after every run: each block is either referenced, free, or cached-and-evictable,
never two of these and never none.

% =====================================================================
\section{Simulating serving policies at scale}\label{sec:serving:sim}

Our engine runs on a CPU with toy models; to study policies at production scale we replace the model
with a \emph{cost model} and keep everything else. \ident{llmsys.serving.simulator} feeds a Poisson
request trace through the \emph{same} \ident{Scheduler} and \ident{BlockManager} classes, and advances a
virtual clock by
\begin{equation}
  t_{\text{step}} = t_0 + \frac{1}{\epsilon}\max\!\left(\frac{2N\,n_{\text{tok}} + \text{attn}}{\Pi},\ \frac{Nb + \kappa\,C_{\text{ctx}}}{\beta}\right)
\end{equation}
where $n_{\text{tok}}$ is the number of tokens in the step, $C_{\text{ctx}}$ the total context read by all
scheduled sequences, $t_0$ a fixed per-step overhead (\SI{2}{ms}) and $\epsilon=0.6$ an achieved
fraction of peak. \Cref{fig:serving:sim} compares three policies for Qwen3-8B on an H100 with a
\SI{20}{GB} KV pool (8,477 blocks of 16 tokens), prompts of 64--512 and outputs of 16--256 tokens:
static batching ($B=32$), continuous batching with paged memory, and continuous batching that
\emph{reserves} 4,096 tokens per request at admission (the pre-paging design).

\begin{figure}[ht]
\centering
\begin{tikzpicture}
\begin{groupplot}[group style={group size=2 by 1,horizontal sep=1.7cm},width=0.49\linewidth,height=5.8cm,
  grid=major,grid style={black!10},tick label style={font=\scriptsize},label style={font=\small},
  xlabel={arrival rate (requests/s)},legend style={font=\scriptsize,at={(0.03,0.97)},anchor=north west}]
\nextgroupplot[ylabel={throughput (output tokens/s)},title={\small throughput},ymin=0]
\addplot[thick,mark=*,mark size=1.3pt,warnred] coordinates {(2,266.1) (4,523) (8,1008) (12,1218) (16,1226) (20,1227) (24,1227) (28,1227) (32,1227) (40,1226)};
\addlegendentry{static, $B=32$}
\addplot[thick,mark=square*,mark size=1.3pt,sysblue] coordinates {(2,267) (4,529.7) (8,1042) (12,1538) (16,2017) (20,2482) (24,2931) (28,3366) (32,3786) (40,4581)};
\addlegendentry{continuous, paged}
\addplot[thick,mark=triangle*,mark size=1.6pt,napkin] coordinates {(2,267) (4,529.7) (8,1042) (12,1538) (16,2017) (20,2482) (24,2625) (28,2642) (32,2657) (40,2670)};
\addlegendentry{continuous, reserve 4k}
\nextgroupplot[ymode=log,ylabel={mean TTFT (s)},title={\small time to first token}]
\addplot[thick,mark=*,mark size=1.3pt,warnred] coordinates {(2,1.201) (4,1.486) (8,1.966) (12,8.162) (16,13.74) (20,17.03) (24,19.53) (28,21.31) (32,22.57) (40,24.36)};
\addplot[thick,mark=square*,mark size=1.3pt,sysblue] coordinates {(2,0.0161) (4,0.01695) (8,0.01748) (12,0.01802) (16,0.01877) (20,0.01966) (24,0.02025) (28,0.02107) (32,0.02162) (40,0.0238)};
\addplot[thick,mark=triangle*,mark size=1.6pt,napkin] coordinates {(2,0.0161) (4,0.01695) (8,0.01748) (12,0.01802) (16,0.02852) (20,0.2636) (24,2.033) (28,3.612) (32,4.775) (40,6.481)};
\end{groupplot}
\end{tikzpicture}
\caption{Simulated (not measured) serving of Qwen3-8B on one H100 under the roofline cost model, 600
Poisson requests per point. Static batching saturates at about 1,200 tokens/s with TTFTs of
tens of seconds. Reservation-based continuous batching caps the running batch at about 30 sequences
($20\,\text{GB}/(4096\times\SI{144}{KiB})$) and its TTFT explodes beyond 20 requests/s. Paged continuous
batching keeps TTFT near \SI{20}{ms} while throughput keeps growing (its mean running batch reaches 59).}
\label{fig:serving:sim}
\end{figure}

Three lessons are visible. First, static batching's TTFT is dominated by waiting for the previous
batch's longest request: seconds even at light load, and it collapses as soon as the arrival rate
exceeds its fixed capacity. Second, continuous batching with reservation behaves identically to paging
until KV memory, not compute, limits the batch, then saturates at a throughput set by
$\text{memory}/\text{reservation}$. Third, with paging the batch grows with load (the mean running batch
climbs from 3 to 59 sequences in this sweep) and TPOT rises only gently (\SI{10.3}{ms} to
\SI{13.3}{ms}), which is the roofline story of \cref{tab:inference:decode} playing out in a queueing
system. Lab~19's \code{bench.py} reproduces this figure and lets you vary block size, budget,
memory, and the prefix-sharing pattern.

\begin{pitfall}[Simulators are only as good as their cost model]
Our cost model ignores attention-kernel efficiency at long contexts, sampling cost, CPU scheduling
overhead, and the interference between prefill and decode inside one step beyond the max() term. Use
it to compare policies and to find regimes worth measuring, then validate the conclusions on the real
engine and hardware.
\end{pitfall}

% =====================================================================
\section{Summary}\label{sec:serving:summary}
\begin{itemize}
  \item Batching amortizes weight reads across rows; it raises throughput and costs memory and some
    latency. It does not help per-row work (KV reads, padding, activations).
  \item Batched generation left-pads prompts and masks pad keys; RoPE makes left padding exact,
    absolute position embeddings need per-row position ids. Training pads on the right.
  \item Evaluation and distillation batch different prompts (padding); self-consistency and GRPO batch
    samples of one prompt (no padding, shareable prefix).
  \item Serving engines schedule at iteration granularity (continuous batching), flatten ragged batches
    instead of padding, manage KV memory in pages with reference counts, share prefixes through hash chains
    or radix trees, bound step cost with chunked prefill, and preempt by recomputation or swapping.
  \item Disaggregating prefill and decode decouples TTFT from TPOT at the cost of moving KV caches, which
    is affordable on fast interconnects.
  \item \ident{llmsys.serving} implements all of this in a few hundred lines and is token-for-token
    identical to the reference generator.
\end{itemize}

% =====================================================================
\section{Exercises}\label{sec:serving:exercises}

\begin{exercise}[Per-row positions]\label{ex:serving:positions}
Extend \ident{GPTModel.forward} (in your own subclass) to accept \code{position\_ids} of shape $(B,T)$.
Show that \ident{generate_batch} with left padding then reproduces single-sequence GPT-2 outputs exactly,
and quantify how different the outputs are without the fix.
\end{exercise}

\begin{exercise}[Dropping finished rows]
Implement strategy (3) of \cref{fig:serving:eos}: after each step, remove rows that emitted EOS from the
input, the validity mask and every layer's KV cache. Measure the speed-up over strategy (2) for output
lengths drawn from a geometric distribution with mean 200 and maximum 2,000.
\end{exercise}

\begin{exercise}[Padding efficiency]
Derive the expected output-padding efficiency $\eta$ of a static batch of size $B$ when output lengths are
i.i.d.\ uniform on $[a,b]$. Evaluate for $B\in\{1,8,64\}$ with $a=50$, $b=1500$, and compare with
continuous batching.
\end{exercise}

\begin{exercise}[Block size]
Using the simulator with the paged policy, sweep the block size over $\{1,4,16,64,256\}$ at a fixed KV
budget. Plot throughput and the internal fragmentation. Then argue, from how a GPU kernel would gather
blocks, why real engines do not use block size 1.
\end{exercise}

\begin{exercise}[Copy-on-write fork]
Add \code{fork(rid, n)} to the engine: it creates $n$ child sequences that share the parent's blocks
after its prefill, and implements copy-on-write for the partially filled last block. Verify with
temperature sampling and fixed seeds that each child generates the same tokens as an independent request
with the same seed, and report the memory saved for $n=16$.
\end{exercise}

\begin{exercise}[Swap vs.\ recompute]
For Qwen3-8B on an H100 with PCIe~5.0 (about \SI{50}{GB/s} effective), compare the per-token cost of
swapping a preempted sequence's KV cache out and back in with the cost of recomputing it. Include the
quadratic attention term of recomputation and find the sequence length at which the cheaper option
changes, if any. Repeat for the same model with full multi-head attention (32 KV heads) and for an fp8
KV cache.
\end{exercise}

\begin{exercise}[Budget trade-off]
In the simulator, sweep the token budget $\tau\in\{64,256,1024,4096\}$ at a load near saturation and
plot mean TTFT against mean TPOT. Explain the shape of the trade-off curve in terms of
\cref{fig:serving:chunked}.
\end{exercise}

\begin{exercise}[Multi-turn chat]
Write a trace generator for multi-turn conversations (each turn appends 50--200 user tokens to the
previous conversation including the model's reply). Compare prefill tokens and TTFT with and without
prefix caching as the number of turns grows. Why does the benefit grow with conversation length?
\end{exercise}

\begin{exercise}[Disaggregation sizing]
A workload has mean prompt 3,000 tokens and mean output 300 tokens on Qwen3-8B. Using the roofline
bounds, estimate the ratio of prefill GPUs to decode GPUs that balances the two pools at a target TPOT
of \SI{30}{ms}. What changes for a reasoning workload with prompt 300 and output 8,000?
\end{exercise}

\begin{exercise}[Fairness]
Implement a Virtual-Token-Counter-style scheduler: track tokens served per tenant and always admit from
the tenant with the smallest counter. Show with the simulator that a tenant sending 10$\times$ more
requests no longer degrades the TTFT of a light tenant.
\end{exercise}

% =====================================================================
\section{Lab 19: A paged, continuous-batching engine}\label{sec:serving:lab}
\begin{labbox}{19}{A paged, continuous-batching engine}
\labdir{lab19_engine}
\labsection{Goals}
Implement the four components that make \ident{llmsys.serving} work, plug them into \ident{LLMEngine},
and show that the engine is token-for-token identical to single-sequence generation under chunked
prefill, prefix sharing and preemption. Then use the simulator to study policies at scale.

\labsection{Setup}
\begin{lstlisting}[style=shell]
cd code
python3 -m pytest -q labs/lab19_engine        # grades starter.py (CPU, < 1 min)
python3 labs/lab19_engine/bench.py --help     # simulator experiments
\end{lstlisting}

\labsection{Tasks}
\begin{enumerate}
  \item \textbf{BlockManager} (\ident{_take}, \ident{allocate}, \ident{release}, \ident{lookup_prefix},
    \ident{register}): free list, reference counts, evictable LRU set, chained-hash prefix lookup that
    always leaves one token to compute.
  \item \textbf{Paged addressing}: \ident{slot_mapping} (\cref{eq:serving:slot}) and
    \ident{paged_attention} with GQA head expansion and a causal offset for chunked prefill. The test
    compares against \code{scaled\_dot\_product\_attention} on a dense copy.
  \item \textbf{Scheduler.schedule}: running sequences first under the token budget, preempt the youngest
    on block shortage, then FCFS admission with prefix lookup. \ident{build_engine} wires your classes in.
  \item \textbf{Fragmentation}: \ident{waste_reserved}, \ident{waste_paged}, \ident{blocks_for_budget}.
\end{enumerate}

\labsection{Measurements}
\begin{itemize}
  \item \code{bench.py --exp policies}: reproduce \cref{fig:serving:sim} (static vs.\ continuous paged vs.\
    continuous reserved) and report the knee of each curve.
  \item \code{bench.py --exp prefix}: requests sharing $k$ system prompts of length $P$; plot prefill
    tokens and TTFT with and without prefix caching for $k\in\{1,4,64\}$.
  \item \code{bench.py --exp budget}: TTFT vs.\ TPOT as the token budget varies.
  \item \code{bench.py --exp engine}: run the real engine on a small random model on your machine; report
    steps, mean batch size, prefix hit rate and preemptions for a burst of 32 requests.
\end{itemize}

\labsection{Deliverables}
Passing tests; the four plots; a one-page design note explaining (i) why the engine's greedy outputs
are identical to the reference even under preemption, and (ii) which component you would move to a GPU
kernel first and why.

\labsection{Stretch goals}
\begin{itemize}
  \item Copy-on-write \code{fork} for parallel sampling (exercise above) and its use for self-consistency.
  \item A Triton paged-attention decode kernel (one program per sequence and head, online softmax over
    blocks) that replaces \ident{paged_attention} on GPU.
  \item Swap-based preemption to a CPU block pool, with a policy that chooses swap vs.\ recompute by
    sequence length.
\end{itemize}
\end{labbox}
